%%%%%%%%%%%%%%%%%%%%%%%%%%%%%%%%%%%%%%%%%%%%%%%%%%%%%%%%%%%%%%%
% ABSTRACT AND INTRO UPDATES FOR PHASE 0 INTEGRATION
% Copy-paste these snippets into your paper
%%%%%%%%%%%%%%%%%%%%%%%%%%%%%%%%%%%%%%%%%%%%%%%%%%%%%%%%%%%%%%%

%%%%%%%%%%%%%%%%%%%%%%%%%%%%%%%%%%%%%%%%%%%%%%%%%%%%%%%%%%%%%%%
% 1. ABSTRACT UPDATE
%%%%%%%%%%%%%%%%%%%%%%%%%%%%%%%%%%%%%%%%%%%%%%%%%%%%%%%%%%%%%%%

% OPTION A: Phase 0 only (use NOW before Overcooked)
% Insert as second-to-last sentence of abstract:

In addition to introducing the O/R Index, we provide a theoretical analysis
characterizing its range and extremal cases, including monotonicity under
noise mixing, and validate these properties in a simple 2$\times$2 coordination
game alongside large-scale multi-agent navigation experiments.

% OPTION B: Phase 0 + Overcooked (use AFTER Phase 1 complete)
% Swap to this version once Overcooked results are integrated:

% In addition to introducing the O/R Index, we provide a theoretical analysis
% characterizing its range and extremal cases, including monotonicity under
% noise mixing, and validate these properties empirically across a 2$\times$2
% coordination game, large-scale multi-agent navigation, and Overcooked-AI
% sequential cooperation tasks.


%%%%%%%%%%%%%%%%%%%%%%%%%%%%%%%%%%%%%%%%%%%%%%%%%%%%%%%%%%%%%%%
% 2. INTRODUCTION: CONTRIBUTIONS BULLET
%%%%%%%%%%%%%%%%%%%%%%%%%%%%%%%%%%%%%%%%%%%%%%%%%%%%%%%%%%%%%%%

% OPTION A: Standalone theory bullet (recommended)
% Add to your "Our contributions are:" list:

\item \textbf{Theoretical analysis and toy validation.}
We characterize the O/R Index mathematically, showing that it takes values
in $[-1,\infty)$ with $-1$ for deterministic observation-based policies and
$0$ for observation-independent policies, and prove that adding private
observation-independent noise monotonically increases O/R. We validate these
properties in a 2$\times$2 coordination game, where O/R moves smoothly from
$-1$ to $\approx 0$ as noise increases.

% OPTION B: Merged with existing O/R introduction bullet
% If you already have "we introduce the O/R Index," merge to:

% \item \textbf{A theoretically grounded coordination metric.}
% We formalize the O/R Index as a normalized residual variance of the
% action distribution under partial observability, prove basic properties
% (range, extremal cases, and monotonicity under noise mixing), and
% illustrate these behaviors in a 2$\times$2 coordination game where
% O/R increases smoothly from $-1$ (deterministic policies) to $0$
% (observation-independent random policies).


%%%%%%%%%%%%%%%%%%%%%%%%%%%%%%%%%%%%%%%%%%%%%%%%%%%%%%%%%%%%%%%
% 3. INTRODUCTION: CLOSING PARAGRAPH (END OF INTRO)
%%%%%%%%%%%%%%%%%%%%%%%%%%%%%%%%%%%%%%%%%%%%%%%%%%%%%%%%%%%%%%%

% OPTION A: Phase 0 + anticipating Overcooked (use NOW)
% This version is valid now and anticipates Section 5.X for Overcooked:

Taken together, these results suggest that the O/R Index is a simple yet
informative diagnostic of coordination quality in multi-agent RL. In
Section~\ref{sec:theory_or} we introduce the O/R Index and analyze its
theoretical properties, including its range, extremal cases, and behavior
under noise mixing, with a concrete 2$\times$2 coordination game as a
running example. We then show in Section~\ref{sec:experiments_nav} that
O/R strongly predicts coordination success in large-scale cooperative
navigation, admits early prediction from partial training trajectories,
and is robust to observation noise and communication dropout. Finally,
we validate that the same pattern extends beyond spatial navigation by
applying O/R to Overcooked-AI sequential coordination tasks
(Section~\ref{sec:experiments_overcooked}), and discuss its practical use
as a sample-efficient diagnostic for MARL pipelines.

% OPTION B: Phase 0 only (conservative, no Overcooked mention)
% Use this if you want to submit before Overcooked runs:

% Taken together, these results suggest that the O/R Index is a simple yet
% informative diagnostic of coordination quality in multi-agent RL. In
% Section~\ref{sec:theory_or} we introduce the O/R Index and analyze its
% theoretical properties, including its range, extremal cases, and behavior
% under noise mixing, with a concrete 2$\times$2 coordination game as a
% running example. We then show in Section~\ref{sec:experiments_nav} that
% O/R strongly predicts coordination success in large-scale cooperative
% navigation, admits early prediction from partial training trajectories,
% and is robust to observation noise and communication dropout. Finally,
% we discuss how these properties make O/R a practical, sample-efficient
% diagnostic for MARL pipelines, and outline how the metric can be extended
% to more complex settings beyond spatial navigation.


%%%%%%%%%%%%%%%%%%%%%%%%%%%%%%%%%%%%%%%%%%%%%%%%%%%%%%%%%%%%%%%
% NOTES ON INTEGRATION
%%%%%%%%%%%%%%%%%%%%%%%%%%%%%%%%%%%%%%%%%%%%%%%%%%%%%%%%%%%%%%%

% Abstract: Insert as second-to-last sentence (before "Broader impacts" etc.)
% Contributions bullet: Add to list in Section 1 (Introduction)
% Closing paragraph: Replace current ending of Introduction section

% After pasting these:
% 1. Update section references (\ref{sec:theory_or}, \ref{sec:experiments_nav}, etc.)
% 2. Verify they compile without LaTeX errors
% 3. Check that cross-references resolve correctly

%%%%%%%%%%%%%%%%%%%%%%%%%%%%%%%%%%%%%%%%%%%%%%%%%%%%%%%%%%%%%%%
% END OF UPDATES
%%%%%%%%%%%%%%%%%%%%%%%%%%%%%%%%%%%%%%%%%%%%%%%%%%%%%%%%%%%%%%%
